\documentclass[../../../include/open-logic-section]{subfiles}

\begin{document}

\olfileid{sth}{replacement}{intro}
\olsection{مقدمه}

جایگزینی طرحوارهٔ اصل موضوعی است که تفاوت میان~$\ZF$ و~$\Z$ را رقم
می‌زند. در سراسر~\crefrange{sth:ordinals::chap}{sth:spine::chap} بدون
درنگ از آن بهره گرفتیم. در این فصل سرانجام این پرسش را بررسی می‌کنیم:
آیا جایگزینی موجه است؟

برای روشن‌کردن پرسش، شایسته است توجه کنیم که جایگزینی واقعاً بسیار
\emph{نیرومند} است. در این فصل (و بار دیگر در
\olref[sth][card-arithmetic][fix]{sec}) تصویری از اندازهٔ این قدرت به دست
خواهیم آورد. اما همین امر نشان می‌دهد که واقعاً به توجیه نیاز داریم.

دو نوع توجیه را بررسی خواهیم کرد. به‌بیان تقریبی، توجیه \emph{بیرونی}
کوششی است برای توجیه یک اصل موضوع بر پایهٔ ثمراتش؛ و توجیه
\emph{درونی} کوششی است برای نشان‌دادن اینکه خودِ مفاهیم ریاضیِ مورد بحث
آن اصل را موجه می‌سازند. در طول این فصل معنای این تمایز روشن‌تر خواهد
شد، اما آنچه می‌گوییم تنها نوک کوه یخ است. برای مطالب بیشتر، به‌ویژه
بنگرید به~\citeauthor{Maddy1988a} (\citeyear{Maddy1988a} و
\citeyear{Maddy1988b}).

\end{document}
